\section{Integral charts at the actual arithmetic points}
\label{sec:actual-good3}

The local theta arguments require a formal chart of the same integral
semiabelian model that carries the prescribed tangent frames and theta
generator.  We now construct that chart and its degeneration.  In
particular, the hypothesis called \texttt{Good3} is supplied at every
original integral framed point after an unramified local extension.

Fix the original host, coefficient field, moduli space and primitive
line.  Let $R$ be a complete discrete valuation ring of characteristic
zero, with fraction field $K$, finite residue field $\kappa$, and the
given embedding $c.F\hookrightarrow K$.  Let $y$ be an original
\texttt{IntFr} point.  Thus its data include an actual semiabelian
model $\mathcal G/R$ of the generic moduli object, its original
rigidified cubical line $\mathscr L$, integral bases
$u_{\tau i}$ of the separate $\tau$-eigenspaces of its Lie algebra,
and a section $\theta$ generating $H^0(\mathcal G,\mathscr L)$.
Their generic identifications with $y$ are part of the data.
Write $w=w_0\in\mathcal O_E$, where $w+\bar w=1$.

\begin{theorem}[Actual integral-chart supplier]
\label{thm:actual-good3}
There is a finite unramified extension of complete discrete valuation
rings $R\subset R'$, with fraction fields $K\subset K'$, and a
base-changed integral framed point $y'$ with the exact original
\texttt{FrMap}.  The ramification index is one and
\[
                      v_{R'}(a)=v_R(a)\qquad(a\in K^\times).
\]
For every integral model in the original \texttt{IntData} of $y'$,
there is an original \texttt{PELUnif} with \texttt{TorsionIso} for
every integer $k\geq1$.  There is an actual
\texttt{IntegralChartA3} $C$ at $y'$ whose frames obey all truncation
identities.  Consequently $C$ satisfies the full original
\texttt{Good3} condition.

Every integral section of every power of the original primitive line
has an integral formal germ in this chart.  The original eigenspace
bases and theta generator are pulled back exactly; no lattice index
or scalar rescaling is removed.  Harmonic theta values computed by
any original \texttt{IsNab} rule commute with this base change.
\end{theorem}

\begin{proof}
\emph{Splitting over an unramified extension.}
Let $T_\kappa$ be the torus in the closed fibre of $\mathcal G$.
The continuous Galois action on its free character lattice has finite
image, since the source is compact and the lattice is discrete.
It therefore becomes trivial over a finite extension $\kappa'/\kappa$.
The equivalence between finite \'etale algebras over a henselian local
ring and over its residue field lifts $\kappa'$ to a finite \'etale
local $R$-algebra $R'$\cite{good3-henselian}.
This ring is a complete discrete valuation ring, its uniformizer is
the original uniformizer, and its residue field is finite.
In particular the valuation scaling is exactly one.  The closed torus
and its Raynaud torus over $R'$ are split.

The prescribed framed data survive this extension without a change of
lattice.  To check the eigenspaces, take an integral basis
$b_1,\ldots,b_d$ of $\mathcal O_E$.  The integral $\tau$-eigenspace is
\[
 U_\tau=\ker\left(\operatorname{Lie}(\mathcal G/R)
       \longrightarrow\operatorname{Lie}(\mathcal G/R)^d,\quad
       v\longmapsto(([b_i]-\tau(b_i))v)_i\right).
\]
Every $\tau(b_i)$ belongs to $R$ by integral closedness.  Flat base
change commutes with this kernel, so the specified $u_{\tau i}$ remain
its exact basis over $R'$.  The assertion does not identify
$\bigoplus_\tau U_\tau$ with the full integral Lie algebra.
The scheme $\mathcal G$ is quasi-compact and separated; flat base
change for sections therefore gives
\[
 H^0(\mathcal G_{R'},\mathscr L_{R'})
       =H^0(\mathcal G,\mathscr L)\otimes_R R'=R'\theta
       \quad\text{\cite{good3-flat-sections}}.
\]
All rigidified line comparisons are pulled back, so this is the same
generic theta generator in the required \texttt{FrMap}.  A finite-type
test realization of $y$ remains a test realization after composition
with $K\hookrightarrow K'$.  This proves every clause of
\texttt{IntFr} and \texttt{IsExt}, with $e=1$.

\emph{Uniformization of every integral model.}
After its splitting paragraph, the proof of
Theorem~\ref{thm:dvr-original-pel-unif} uses only a split closed torus,
the actual original cubical line and its primitive identity, and a
principally polarized abelian generic fibre.  Its level assumption
was used solely to obtain splitting over the original base.
Apply the remainder of that proof to $\mathcal G_{R'}$.
In particular, ampleness is supplied before the ample degeneration
equivalence; the degree identity
\[
       1=\deg\lambda=[X:\phi(Y)]^2\deg\lambda_B
\]
identifies the period and character lattices by the original
$\phi$, and the original primitive identity fixes the character
ambiguity of the coefficient line.  The exact all-thickening
expansions, their common denominators, and the original torsion
translation formula follow with every integer $k$, including the
residue characteristic.  No assertion that the entire torsion kernel
is finite flat is required.

If $\mathcal H/R'$ is any other \texttt{IntData} model of $y'$,
the two given generic identifications give an isomorphism
$\mathcal H_{K'}\simeq\mathcal G_{K'}$.
Homomorphisms of semiabelian schemes over a normal base extend
uniquely from a dense open.  Extend this isomorphism and its inverse;
their compositions are the identity by the same uniqueness.
The actual generic cubical comparison of the primitive lines also
extends uniquely\cite[Propositions 3.3.1.5 and 3.3.2.2]{good3-lan}.
This uses that comparison, not zero rigidification alone.
It carries the original action, adjoint identifications and primitive
identity to those of $\mathcal H$.  Hence its closed torus is split
and the same uniformization construction applies to \emph{every}
such model.  This proves the universal model quantifier in the
original extension supplier.

\emph{One integral chart and one compatible sequence of frames.}
Choose a basis of the full free module
$\operatorname{Lie}(\mathcal G_{R'}/R')$, independently of the
eigenspace bases.  Smoothness at the identity gives formal coordinates
with this tangent basis:
\[
                  \widehat{\mathcal G}_{0}
                         \simeq\operatorname{Spf}R'[[x_1,\ldots,x_g]]
                         \quad\text{\cite{good3-formal-coordinates}}.
\]
The formal group law and the $\mathcal O_E$-action therefore have
coefficients in $R'$.  Record each prescribed $u_{\tau i}$ by its
integral coordinates in this full basis.  This is precisely the
original \texttt{uvecT}; its matrix need not be unimodular.
The coordinate map gives all jet maps $x_m$ by restriction.

On the formal product, the line $\delta\mathscr L$ is free over
the local ring $R'[[x,z]]$.  Choose a generator and normalize it on
both axes: if its axis restrictions are $a(x)$ and $b(z)$ relative
to the prescribed rigidifications, multiply by
$a(0)/(a(x)b(z))$.  Compatibility at the origin makes this generator
equal to one on each axis.  Its restrictions to the jets give a
single compatible system $e_m$.  Thus \texttt{FrTrunc} holds as an
actual equality for every $m$, not merely after separate choices of
frames on each thickening.

The ratios of the original two addition maps, symmetry, and adjoint
maps in this frame are integral unit series
$\phi_1,\phi_2,\mathrm{sy},\mathrm{ad}(a)$.  Their identities are
the literal identities of the original cubical biextension, so they
give the required integral formal-biextension data.  Restrict the
frame along $(x,[w]x)$ and apply the \emph{given} primitive identity.
This is the original frame of $\mathscr L$, with value one at zero.
Every integral section of $\mathscr L^k$ consequently has an integral
power-series germ, for every $k\geq0$.
All these constructions commute with passage to the generic fibre.
It remains to supply the canonical gauge required by
\texttt{ChartAt}; its existence is proved next over $K'$.
\end{proof}

\paragraph{An explicit canonical gauge.}
We give the construction because integral chart coordinates alone
do not supply this part of the original connection interface.
The logarithm of a commutative formal group over a characteristic-zero
field gives additive coordinates with tangent map the identity.
For completeness, the matrix of invariant differentials is the
inverse of the derivative of the group law in its second variable
at zero.  Associativity gives invariance and commutativity makes these
forms closed.  Integrating their homogeneous pieces by the formal
Euler homotopy gives the logarithm.  All divisions in this paragraph
take place in $K'$, not in the integral chart ring.

Use these additive coordinates temporarily.  A normalized symmetric
formal additive two-cocycle $c(x,x')$ is a coboundary.  Indeed the
cocycle identity makes
$-\partial_{x'}c(x,0)\,dx$ closed; its formal integral $H$ with zero
constant and linear terms satisfies
\[
                  c(x,x')=H(x)+H(x')-H(x+x').
\]
The integration is linear in the cocycle and works coefficientwise
with additional formal parameters.  Applying it first to
$\log\phi_1(x,x',z)$ kills the first addition law by a gauge $H_1$.
The biextension interchange law then makes the logarithm of the
remaining second addition cocycle additive, hence linear, in $x$.
Apply the same homotopy in its two $z$ variables, coefficientwise in
$x$, to obtain $H_2$ linear in $x$.  It kills the second addition
law and leaves the first one unchanged.  Both gauges vanish on the
axes.

After these changes, $\delta(a)=\log\mathrm{ad}(a)$ is bilinear.
Let $M$ be the vector space of bilinear forms on the formal Lie
algebra, with commuting actions
\[
        (l_a B)(x,z)=B(ax,z),\qquad
        (r_a B)(x,z)=B(x,\bar a z).
\]
The original adjoint identities give
\[
 \delta(a+b)=\delta(a)+\delta(b),\qquad
 \delta(ab)=l_b\delta(a)+r_a\delta(b).
\]
They extend from $\mathcal O_E$ to $E$ by rational linearity.
Choose a trace-dual pair of bases of the separable algebra
$E/\mathbb Q$ and write its separability idempotent as
$\sum_i u_i\otimes v_i$.  Thus $\sum_i u_iv_i=1$ and
$\sum_i au_i\otimes v_i=\sum_i u_i\otimes v_i a$.
Set
\[
                         H_3=\sum_i l_{u_i}\delta(v_i).
\]
Commutativity of $E$ gives also
$\delta(ab)=l_a\delta(b)+r_b\delta(a)$.
Applying the idempotent identity to the map
$u\otimes v\mapsto l_u\delta(v)$ yields
$l_aH_3=\delta(a)+r_aH_3$.  Therefore
$\delta(a)=l_aH_3-r_aH_3$, the exact adjoint coboundary.
This bilinear gauge preserves both addition laws.

The remaining $s(x,z)=\log\mathrm{sy}(x,z)$ is bilinear,
skew-symmetric and balanced for the adjoint action.  The gauge
$H_4=-s/2$ preserves both addition laws and the adjoint identity,
and satisfies $H_4(z,x)-H_4(x,z)=s(x,z)$.
Transporting the sum of the four gauges back to the original
coordinates gives a series $h$, zero on both axes, satisfying
exactly
\begin{align*}
 \phi_1(x,x',z)&=
   \exp\bigl(h(x,z)+h(x',z)-h(F(x,x'),z)\bigr),\\
 \phi_2(x,z,z')&=
   \exp\bigl(h(x,z)+h(x,z')-h(x,F(z,z'))\bigr),\\
 \mathrm{sy}(x,z)&=\exp\bigl(h(z,x)-h(x,z)\bigr),\\
 \mathrm{ad}(a;x,z)&=
   \exp\bigl(h([a]x,z)-h(x,[\bar a]z)\bigr).
\end{align*}
These are the source's four \texttt{IsCanonGauge} identities, with
their original signs\cite{good3-source}.  Together with the integral
chart above they finish every clause of \texttt{IntegralChartA3}.
Its compatible frames and the proved \texttt{PELUnif/TorsionIso}
are exactly the two clauses of \texttt{Good3}.

\paragraph{The original harmonic values and local denominators.}
Two canonical gauges for the same biextension differ, in logarithm
coordinates, by a balanced symmetric bilinear form $B(x,z)$.
The difference on the primitive graph is
\[
 B(x,wx)=\tfrac12 B(x,x),
\]
because $B$ is symmetric, adjoint-balanced, and $w+\bar w=1$.
A change of formal coordinates or of biextension frame first
transports the actual section germ and gauge; their remaining
ambiguity is precisely such a bilinear form.
For tangent eigenvectors, balancing implies
$B(u_{\tau i},u_{\nu j})=0$ unless $\nu=\bar\tau$.
Consequently the effect on a column product of ordinary derivatives
is the finite contraction operator
\[
 \exp\!\left(\sum_{\tau,i,j} b_{\tau ij}
       \sum_{b=1}^{n}
       \frac{\partial^2}
        {\partial X_{\tau i b}\partial X_{\bar\tau j b}}\right)
\]
applied to the polynomial prescribing that value.  This identity
follows by multiplying the generating series by the displayed
quadratic exponential; it is a finite identity in each degree.
Every positive-order contraction vanishes for the original harmonic
polynomials.  The original \texttt{IsNab} value may therefore be
computed in the integral chart just constructed.  This argument
asserts no gauge independence for arbitrary nonharmonic derivatives.

Choose the chart and gauge over $K$ and pull them to $K'$ when
comparing the two generic values.  The same harmonic invariance
compares them with any chart chosen by the original \texttt{IsNab}
rule at either point.  Hence $\iota(H_y)=H_{y'}$ with the exact
original section, tangent vectors and polynomial coefficients.
Normalized valuations are unchanged because $e=1$.

In particular Theorem~\ref{thm:discriminant-harmonic-clearing}
now applies to the \emph{actual} chart and every actual integral
section in it.  If $P$ has total degree $d$ and
$\Delta=|\operatorname{disc}E|$, it gives
\[
                         \Delta^d H(P;f)\in R'.
\]
When the section data and their value descend to $K$, this also
belongs to $R$, since $K\cap R'=R$.
Thus passage to the splitting extension introduces no additional
valuation loss.  The torus rank, the eigenspace lattice defect and
the existing discriminant exponent are all retained.

This supplies the local chart and uniformization gate used by the
actual primary-trace and reached-section arguments.  It does not
bound the degree of the local unramified extension uniformly, replace
the global coefficient-field bounds, or by itself construct the
entire arithmetic cover.  The original theta-series identification,
determinant-ideal scalars and global norm estimates retain their own
separate obligations.
